% !PRINTED 398
\weberpage{398}
%<<CONT
real double tangents $z_3, z_4$, so that the complex 4) continues
as follows:
\[
4)\quad x_1z_1,\ x_2z_2,\ x_3z_3,\ x_4z_4;
\]
$z_3$ and $z_4$ do not occur in 2), nor in 1) either, and
must therefore be contained in 3).

Neither can they occur paired in 3), because
$x_1, z_3, z_4$ are azygetic. From this, if $t_3, t_4$ are two
further real double tangents, we conclude that the complex 3)
is composed as follows:
\[
3)\quad x_1y_2,\ x_2y_1,\ z_3t_3,\ z_4t_4.
\]

If one now considers, in place of the complex 4), the complex
\[
5)\quad x_1z_3,\ y_2t_3,
\]
then one can conclude in exactly the same way that there is a
fourth real pair $u_1u_2$ in the complex 2), and with this 5. is proved.

If we gather the result together, we recognise that,
with regard to the reality of the double tangents of a real
curve of the fourth order, only four cases are possible:

1) \gesperrt{Four real syzygetic double tangents.}

2) \gesperrt{Eight real double tangents, namely four pairs
of a Steiner complex.}

3) \gesperrt{Sixteen real double tangents, which form four
real pairs each in a syzygetic complex triple.}

4) \gesperrt{Twenty-eight real double tangents.}

\begin{center}
\S.\,105.\\
\gesperrt{Proof of the existence of the four cases.}
\end{center}

In the preceding section we have at first proved only
that there can be no cases other than 1), 2), 3), 4).
But that these four cases are indeed all possible is now
also easy to show, on the basis of the theorem of \S.\,100, according
to which from seven arbitrarily given straight lines one can
derive, by a \gesperrt{rational way}, a curve of the fourth order for
which the given lines form an \gesperrt{Aronhold} system,
from which all the double tangents can be derived \gesperrt{rationally}.

Among the seven given straight lines imaginary ones may
also occur, and as a condition for the reality of the curve
%CONT>>
